\section{Interpretation and limitations}
\label{sec:discussion}

\subsection{A prediction-targeted smoothing choice}

Normalized smoothness describes the effective action
of the historical PLS smoother. For fixed $(d,L)$,
\begin{equation}
\label{eq:scope_smoothness_edf}
S=\frac{L-\operatorname{tr}\{H(S)\}}{L-d}.
\end{equation}
A value near zero preserves much of the historical
variation; a value near one approaches the
least-squares polynomial null space. When
the trend is to be forecast, the operative
criterion is the accuracy of its continuation,
not just its apparent historical regularity.
The same series may therefore warrant
different smoothness values at different
forecast horizons.

The continuation remains a discrete polynomial
of degree at most $d-1$ for fixed difference
order $d$. Changing $S$ changes the fitted
terminal level and differences and therefore
the coefficients of that future polynomial.
Forecast-CV does not, by itself, select the
polynomial degree. Joint order selection or
a stochastic residual forecast requires
a separately specified validation design.

\subsection{What the compact index does and does not do}

Using $S\in[0,1]$ avoids a finite arbitrary
upper boundary on $\lambda$ and gives the two
limiting smoothers an explicit role in
validation. However, $S(\lambda)$ is one-to-one:
the reparameterization preserves all fitted
trends and interior stationary points.
It is not a new smoothing estimator, it
does not guarantee a unique minimizer,
and it does not itself produce predictive
improvement. Large changes in the
unbounded penalty may be compressed
near $S=1$, so numerical search must
treat boundary regions carefully.

\subsection{Simulation scope and statistical caution}

The main evidence is a factorial Monte Carlo
comparison of smoothness \emph{selection rules}
for one specified PLS forecast-continuation
family. The paired decisions share seeds
and, in some conditions, innovations.
Thus the reported 72,000 outer evaluations
are not independent observations.
The aggregate geometric forecast-RMSFE
ratios do not establish uniform or
distribution-free optimality.

The baseline uses an identity-weight
PLS fit with equally spaced observations,
a fixed training length, fixed difference
order, and deterministic polynomial
continuation. AR(1) and heavy-tailed
disturbances enter the simulation, but
the fitting method does not estimate
their complete covariance structure.
Forecast-targeted smoothing cannot
eliminate all end-of-sample and
model-misspecification issues
\citep{hamilton2018hp}.
In particular, the method can perform
poorly if the true future evolution
is incompatible with its extrapolation
class.

